\section{Constructing the integral decomposition at large primes}

% BEGIN D79G1
\paragraph{D79G1. The Hecke coefficients and full matrix span of the actual complement.}
Fix a native member with its $\mathrm{Facts}(D,\mathrm{nch})$, a
localization $i$ with $\mathrm{HG}$, a good representation family
$R$, $\pi\in Z_i$, and an actual continuous rank-$r$ complement
$\rho$ of $R_\pi$. Write $O$ for the coefficient DVR, $K$ for
its fraction field, $k$ for its residue field, and $T\subset O^X$
for its local finite Hecke algebra. The facts give the coefficients of
$\det(1-z R_x(g))$ as elements of $T$ for every $g$: this is the
Hecke coefficient field of $\mathrm{Conclusion41}$ applied to
that good family. At $\pi$ the specified splitting gives the product
of the characteristic polynomial of $\rho(g)$ and the $t$ prescribed
character factors. The same conclusion in the facts makes the
residual characteristic polynomial of $R_x(g)$ independent of $x$.
Consequently the reduction of $\rho$ has the required residual
Frobenius polynomials: reduce the split product at $\pi$, use this
independence at $x$, and insert the good family's Frobenius identity
there. Its kernel is open by continuity of the integral representation
and discreteness of the residue field. These are precisely the two
conditions defining the residual rank-$r$ representation in
$\mathrm{IsRhoBar}$. The member's $\mathrm{rhoBar}$ field at
$\mathrm{HG}$ therefore makes the scalar extension to $\bar k$
irreducible.

It follows that the $O$-span of $\rho(g)$ is $M_r(O)$.
Over $\bar k$ the span is an algebra, since products are group
matrices. Its Jacobson radical annihilates the simple module
$\bar k^r$, on which the algebra acts faithfully. Hence the algebra
is semisimple; by the finite-dimensional structure theorem over an
algebraically closed field it is a product of full matrix algebras.
A faithful simple module sees exactly one factor, which acts by
all endomorphisms. Full span descends to $k$, since matrix rank
is unchanged under extension of fields. Choose $r^2$ group matrices
giving a basis after reduction. Their determinant after vectorizing
is a unit of $O$, so they form an integral basis as well. This proves
both inputs at the actual complement; it uses no choice of a
simultaneous splitting at the other Hecke points.
% END D79G1

% BEGIN D79G2
\paragraph{D79G2. Determinants of Hecke-linear combinations in residue characteristic.}
Suppose that on a finite set of points the group matrix family $R$
satisfies $\det(1-zR(g))\in T[z]$ for every $g$. Let
$A=\sum_g T R(g)$, a subalgebra of the matrix families.
For any finite $a_j\in A$ all coefficients of
$\det(\sum_j u_j a_j)$ belong to $T$. The same holds for
the determinant of a matrix of arbitrary square blocks from $A$.
Both assertions persist after pulling $R$ to a cover and enlarging
the coefficient algebra to contain the pullback of $T$.

For the first assertion, start with formal noncommuting letters
$X_1,\ldots,X_s$. In the algebra of integral formal series in the
words, $1-\sum X_i$ can be factored to any prescribed degree $N$
as a finite ordered product of factors $1+d_w w$, where $w$ is
a nonempty word and $d_w\in\mathbb Z$. In fact start with the
empty product, and, in increasing degree, correct each coefficient
by multiplying by the factor whose integer is its discrepancy.
Because the constant term is one, a degree-$a$ correction changes
only its chosen coefficient up to that degree. Write the original
linear combination as $\sum f_i R(g_i)$ where each $f_i$ is a
polynomial over $T$ in the variables $u_j$. Substitute
$X_i=z f_i R(g_i)$ in the truncated factorization. The determinant
of each factor is a specialization of the known polynomial for
the group product $g_w$ by the scalar
$-d_w z^{|w|}\prod_{i\text{ in }w} f_i$. Thus through degree $N$
all coefficients of $\det(1-z\sum f_i R(g_i))$ are in $T[u_j]$.
Its top possible degree is the matrix size, and its coefficient
there gives the first assertion up to sign. The same argument
applies to finite polynomial combinations with coefficients in $T$.

For the block assertion, write $B$ for the block matrix and
successively eliminate the first row and column of $1+zB$.
Each pivot block has constant term one, whose inverse as a formal
series is obtained by the geometric series. The Schur complements
are block matrices with constant term the identity and coefficients
in the same algebra. As in G79E11, after truncation the determinant
of each diagonal pivot lies coefficientwise in $T$, so the required
highest coefficient of the product lies there. This proves both
clauses of the native determinant condition on the Hecke span
integrally, including in small residue characteristic.
% END D79G2

% BEGIN D79G3
\paragraph{D79G3. Integral projectors with the prescribed ranks at every point.}
Assume the Hecke coefficient condition and the full integral span
of D79G1, and fix the given $K$-splitting
$R_\pi=\rho\oplus\chi_1\oplus\cdots\oplus\chi_t$.
For any $P\in M_r(O)$ there is $a\in A$ whose blocks at
$\pi$ are $(P,0,\ldots,0)$. Indeed the map
$O[G]\to M_r(O)$ induced by $\rho$ is onto. The two-sided ideal
generated by commutators kills all character representations.
Its image contains $[E_{11},E_{12}]=E_{12}$, since $r\ge2$,
and hence every matrix unit by multiplication on both sides.
Lift $P$ from this ideal and evaluate the same group-ring element
through the family.

For $b\in A$ idempotent at $\pi$, form the limit of $b^{j!}$.
Modulo each uniformizer power the finite matrix semigroup makes
the powers eventually periodic. For all sufficiently large $j$,
$j!$ is in the periodic range and divisible by the period; hence
the residues stabilize to an idempotent (the residues for $j!$
and $2j!$ agree). Completeness gives an idempotent $e$ in the
ambient finite matrix family. The submodule $A$ is closed:
it is a finitely generated submodule of a finite free module
over the compact DVR $O$, hence is the compact image of a finite
power of $O$. Thus $e\in A$ and $e(\pi)=b(\pi)$. Any left
or right annihilator of $b$ still annihilates $e$ by passage to
the limit. Apply this first to $1-a$ for $P=1_r$ to obtain
$e_B$ projecting onto the character summand at $\pi$. Apply it
again to $(1-e_B)a'(1-e_B)$, where $a'$ lifts $E_{11}$ on the
rank block, to obtain $e_1$. Then $e_Be_1=e_1e_B=0$, and $e_1$
acts at $\pi$ as the chosen rank-one projector on the rank block.

These projectors have ranks $t$ and $1$ at every point.
For $f\in T$ the residues of $f(x)$ are independent of $x$:
the surjections $T\to O\to k$ agree on constants and have the
same kernel, the unique maximal ideal of $T$. The coefficients
of the characteristic polynomial of either projector lie in $T$
by D79G2, applied to its linear combination and the identity.
If an integral idempotent has rank $s$ over $K$ its characteristic
polynomial is $U^{n_V-s}(U-1)^s$ already over $O$. Reduction
retains the exact multiplicity of the root $1$, even in positive
characteristic. It is therefore constant throughout $X$.
% END D79G3

% BEGIN D79G4
\paragraph{D79G4. A single tame cyclotomic generator and all its character roots.}
At a prime $\ell>p_*$ of a member past G79E5's uniform cutoff,
$\ell$ avoids $q,d$ and the host's small primes, and $t+1<\ell$.
The native conductor has its prime factors among these primes,
so $\ell$ does not divide it. The host's Frobenius condition
away from the conductor identifies the identity as the only element
acting trivially on the residue field in the decomposition group.
Since $E/\mathbb Q$ is Galois the inertia order is the ramification
index. Thus $E$ is unramified over $\ell$ at every $w\mid\ell$.

More generally assume this unramifiedness and $t+1<\ell$ directly.
Choose any $w\mid\ell$ and its specified place lift to the algebraic
closure. The extension $\mathbb Q_\ell(\zeta_\ell)/\mathbb Q_\ell$
is totally ramified of degree $\ell-1$, by Eisenstein applied to
$((1+U)^\ell-1)/U$. After the unramified extension $E_w$ the same
polynomial is still Eisenstein, and the extension contains all
conjugates. Inertia at the specified lift therefore surjects onto
$(\mathbb Z/\ell)^\times$ under the cyclotomic character: the
whole finite Galois group is inertia, and inertia in the absolute
decomposition group maps onto it. Choose $g_0$ mapping to a
generator. The character values
$\lambda_a=\chi_a(g_0)=\chi_1(g_0)\chi_{\rm cyc}(g_0)^{-(a-1)}$
for $1\le a\le t$ are pairwise distinct modulo the maximal ideal;
their differences are units, since these exponents are distinct
modulo $\ell-1$.

Set $X_0=e_B R(g_0)e_B\in A$. It is supported on the corner
$e_B$, and at $\pi$ it acts by $\lambda_a$ on the $a$th
character and by zero on the complementary block. The polynomial
\[
 P_x(U)=\det(Ue_B(x)-X_0(x)+1-e_B(x))
\]
is monic of degree $t$: over $K$ it is the ordinary characteristic
polynomial on the image of $e_B$ and the complementary block has
determinant one. Its coefficients belong to $T$ by D79G2.
They have the same residue at every $x$ by D79G3. At $\pi$ the
polynomial is $\prod_{a=1}^t(U-\lambda_a)$. Each simple residue
root lifts by Hensel's lemma over the complete DVR $O$ to a unique
root $\beta_a(x)$ of $P_x$, indexed by its residue. Their residues
are distinct, so these are $t$ distinct integral roots and their
product of monic linear factors equals $P_x$. At $\pi$ they
agree with the original $\lambda_a$ in that order.
% END D79G4

% BEGIN D79G5
\paragraph{D79G5. The ordered-root cover and its integral Hecke basis.}
With the roots of D79G4 put $X'=X\times\mathfrak S_t$,
$p(x,\sigma)=x$, $\pi'=(\pi,1)$, and
$r_a(x,\sigma)=\beta_{\sigma(a)}(x)$ for $1\le a\le t$.
Let $T'\subset O^{X'}$ be generated by $p^*T$ and the $r_a$.
Each fiber of $p$ has $t!$ points. The algebra $T'$ is free of
rank $t!$ over $T$ with basis
\[
  r_1^{m_1}\cdots r_t^{m_t},\qquad 0\le m_a<t+1-a.
\]

Indeed form the universal ordered splitting algebra of $P\in T[U]$.
Adjoin the first root by its monic polynomial of degree $t$;
divide by its monic linear factor and adjoin a root of the resulting
monic polynomial of degree $t-1$, continuing down to degree one.
Monic division proves inductively that the indicated monomials are
a basis, with exactly $t!$ elements. Evaluation on all ordered
roots maps this algebra onto $T'$. It is injective: the algebra
is free over $T$ and $T$ is $O$-torsion-free. After extension to
$K$, $T_K=K^X$ by point separation (G79E11). At each $x$, the
Chinese remainder theorem for the $t$ distinct roots gives the
product of copies of $K$ over all possible choices of the first
root, then over the choices of the remaining roots. Thus evaluation
over $K$ is the isomorphism onto $K^{X'}$. Its integral kernel is
zero. This also gives the uniqueness of all Hecke coefficients in
the basis expansion required by the original decomposition.
% END D79G5

% BEGIN D79G6
\paragraph{D79G6. The full integral decomposition, uniformly at the guarded localizations.}
On the cover of D79G5 take the two projectors, $g_0$, $X_0$, roots,
point and splitting of D79G3--D79G5. At $\pi'$ define the block
maps on the span $A'=\sum_g T' p^*R(g)$ by conjugating evaluation
through the specified $K$-splitting. These maps have values in
$M_r(O)$ and $O$ respectively: for $a=\sum f_j p^*R(g_j)$
the blocks are the same sums with $f_j(\pi')$ multiplying
$\rho(g_j)$ or $\chi_c(g_j)$. Their values are independent of
the expression by the splitting. Off the span extend these maps
arbitrarily, as the native record imposes its identities on the
span alone. The rank-block map is onto by full integral span.
The coefficient conditions for $A'$ follow from D79G2 over $T'$.

Finally put
\[
 E_1=\prod_{a=2}^t (X_0-r_a e_B),\qquad
 \delta=\prod_{a=2}^t(r_1-r_a).
\]
At each point of the cover the corner has $t$ distinct roots over
$K$, so it is diagonalizable there. This polynomial vanishes on
the roots other than $r_1$ and acts by $\delta$ on the $r_1$-line;
it also vanishes on the complementary summand because $t\ge2$.
Thus $E_1$ has rank one and $E_1^2=\delta E_1$ over $K$, hence
the latter identity holds over $O$. At $\pi'$ its image is the
prescribed first-character line and $\delta(\pi')$ is a unit by
D79G4. This verifies the remaining fields of the native
orthogonal-idempotent decomposition, including its determinant
conditions, its ordered-root basis, its integral matrices and
the order of its chosen inertia element. It yields the full original
existence theorem whenever its characteristic-polynomial, residual
span, splitting, unramifiedness and $t+1<\ell$ hypotheses hold;
residual span suffices by lifting a residual matrix basis as in D79G1.

In the setting of G79E10 the same decomposition is available for
every guarded localization, every native good family, every
$\pi\in Z_i$ and every allowed complement. D79G1 uses the given
member facts, and D79G4 uses the family cutoff of G79E5. Neither
step needs the local cut-span conclusions of G79E13.
% END D79G6

% BEGIN D79G7
\paragraph{D79G7. The off-support and ramified-place conditions from the same facts.}
At every localization and complement of D79G1 the member facts
supply all the away-from-$\ell$ hypotheses of G79E13.
Indeed a nonzero vector fixed by $\chi_1^{-1}R_x$ generates a
$\chi_1$-line of $R_x$. Under $\mathrm{HG}$ the
$\mathrm{ConclusionG6}$ field of the facts and the point identity
put such an $x$ in $Z_i$. This gives the off-support condition.
At a finite place outside $S_i,Q_i,\ell$, the last part of
$\mathrm{Conclusion41}$ gives $R_x(\gamma)=1$ for each inertia
element $\gamma$ and every $x$. On the first-character summand
at $\pi$ this forces $\chi_1(\gamma)=1$ too. Therefore the twisted
family is trivial on this inertia as required.

At a prime $P\in Q_i$, the quadratic relation of G79E8 is the
$Q_i$ clause of the member's $\mathrm{ConclusionG5}$, applied at
$\mathrm{HG}$ to this good family and its inertia. The tame generator
has $\zeta_P\ne1$ by the character argument of G79E2. At $P\in S_i$
the defining prime alternatives are a place above one of the two
host type primes,
or a place above the designated auxiliary prime (with the source's
alternate prime when the residue characteristic equals the first one).
In the first case take the Hecke polynomial supplied by
$\mathrm{ConclusionG5}$ for the given tame generator. Its
annihilation relation holds at every point. The value on the first
character is a unit in the Hecke algebra, whose evaluation at $\pi$
remains a unit. Thus the exact alternative $\mathrm{PolyAt}$ holds.

In the designated auxiliary case $\mathrm{ConclusionG5}$ supplies
the quadratic relation. On this inertia the additional twist is
trivial because $S_i\cap Q_i=\varnothing$, and the base character
is the specified auxiliary character $\delta$ or $\delta^{-1}$
under the coefficient map. Its inertia image has order $15$ in
the defining datum. This order persists under the coefficient map:
if the image of $u^{15}=1$ is $1$ and $u\ne1$, the domain identity
$(u-1)(1+u+\cdots+u^{14})=0$ forces the second factor to vanish,
whose image would give $15=0$ in the characteristic-zero target.
Inversion also preserves the order. The residue characteristic of
the auxiliary place is prime to $15$, since its norm is $1$ modulo
$15$. Hence wild inertia has trivial image in this finite group.
The image of the fixed tame generator therefore has order $15$.
For this primitive root $\zeta_P$, $\Phi_{15}(\zeta_P)=0$ and
$\Phi_{15}(1)=1$, the latter following from
$\Phi_{15}(U)=(U^{10}+U^5+1)/(U^2+U+1)$ over $\mathbb Q$.
Subtracting the two polynomial values shows that $\zeta_P-1$
divides $-1$ in $O$. This supplies the unit in the quadratic
alternative, completing all the ramified-place requirements.
% END D79G7

% BEGIN D79G8
\paragraph{D79G8. Reduction of the whole large-prime field to its local conditions above the residue characteristic.}
Assume global Artin reciprocity, as in the original
$\mathrm{Hypotheses}_2$. Fix any native family and use the cutoff
of G79E5. For a qualified
member retain its given $\mathrm{Facts}(D,\mathrm{nch})$. Suppose
that at each localization in $\mathrm{CongBound\_at}$, for every
allowed $R,\pi,\rho$, and each $w\mid\ell$, its native local
cut generators have the following properties:
the rank-block generators lie in the Bloch--Kato finite part,
the generators on character ratios $\chi_{j+1}/\chi_1$ for
$1\le j<t$ are coboundaries, and the first-character generators
are unramified on inertia. Then $\mathrm{CongBound\_at}(D,L)$ holds.

For each such choice D79G6 provides the orthogonal decomposition and
unit root difference. D79G7 provides the other global hypotheses.
Membership of each rank generator in the finite part gives the
rank-block bound with parameter zero by taking their span. G79E13
applied with the given global Artin reciprocity gives the rational
bound when the native lengths are finite. G79E1 gives the extended
bound for all lengths using the nontrivial tame characters from
D79G7. This proves precisely the quantified member assertion.

Consequently, under J78H1 and this uniform local assertion, G79E2
and G79E10 give $\mathrm{Concl.large}$ on the selected hub. The
remaining native task for this route is to prove those three local
conditions above $\ell$ from the member facts and the stipulated
local comparison theorems. G79E5 gives the local numerical position
data, and L78W13 gives the integral Wach functor and comparison.
L79C1--L79C2 descend all three conditions from a coefficient
enlargement. Thus it suffices to identify the actual representation's
jumps, construct the required shifted lattices there, and prove the
three cohomological conditions on the enlarged local cut generators.
% END D79G8
